\textbf{Local gradient interpretation.}
The regression form changes how critic guidance is processed before it
reaches the actor. The following identity makes that distinction precise.

\begin{proposition}
Fix the sampled states and latents, critic, and normalization scale at
$\theta_{\mathrm{old}}$. Let $\Delta_i=\tilde r_i-\bar r_i$ and
$J_i=\partial r_\theta(s,z_i)/\partial\theta$ at these parameters. Then
\begin{equation}
 \left.\nabla_\theta\mathcal L_{\mathrm{actor}}\right|_{\theta_{\mathrm{old}}}
 =-\frac{2}{KD}\mathbb E_s\!\left[\sum_i J_i^\top\Delta_i\right].
 \label{eq:local-gradient}
\end{equation}
\end{proposition}
\textit{Proof.}
The targets are detached. Differentiating \eqref{eq:actor-regression}
therefore gives $2J_i^\top(r_\theta-\tilde r_i)$ per sample; at the
construction point, $r_\theta-\tilde r_i=-\Delta_i$.\hfill$\square$

For $\rho=0$ or $\rho\geq\epsilon$, the restoring field is the negative gradient of
\begin{equation}
 U(r)=\frac{\eta_0}{2}\|r\|_2^2+
       \frac{\eta_{\mathrm{anc}}}{2}
       (\|r\|_2-\sqrt D\rho)_+^2,
 \label{eq:radial-potential}
\end{equation}
with the usual continuous extension at zero. With both clips inactive,
the local actor gradient is $2/D$ times that of the sampled objective with
per-action term $-\eta_Q\bar Q_\phi/q_{\mathrm{scale}}+U(r_\theta)$.
Target regression alone does not supply a
new first-order improvement direction. CAST's explicit field and total-step
clips rescale proposed displacements sample by sample before they pass through
$J_i$. This also explains why a loss-space hinge is a meaningful control,
without implying equivalence of the finite, clipped training procedures.

\textbf{Distribution interpretation and limits.}
Let $P_\theta^s$ and $P_0^s$ be the action distributions induced by the
residual policy and frozen base at state $s$, with finite second moments.
Sharing the latent $z$ defines a valid coupling, so
\begin{equation}
 W_2^2(P_\theta^s,P_0^s)
 \leq \mathbb E_z\|r_\theta(s,z)\|_2^2.
 \label{eq:coupling-bound}
\end{equation}
Here $W_2$ is the 2-Wasserstein distance. The bound follows by evaluating
the transport cost under that coupling rather
than minimizing over all couplings. It connects residual control to
distributional proximity without requiring a policy likelihood.

Neither identity proves that the learned actor stays inside a fixed ball.
The target cap controls $\Delta_i$, while the realized output change after
a parameter step also depends on Jacobian coupling between samples.
Likewise, \eqref{eq:coupling-bound} depends on the \emph{actual} residual,
not merely the requested target step. A nearby policy can still exploit a
wrong critic direction, and a useful policy can lie farther from the base.
Our experiments therefore test stability and improvement empirically.
